\documentclass[11pt]{article}

\usepackage[a4paper,margin=1in]{geometry}
\usepackage{amsmath,amssymb,amsthm}
\usepackage{bm}
\usepackage{mathtools}
\usepackage{hyperref}
\usepackage{graphicx}
\usepackage{xcolor}

\hypersetup{
  colorlinks=true,
  linkcolor=blue!50!black,
  urlcolor=blue!50!black,
  citecolor=blue!50!black
}

\setlength{\parskip}{0.6em}
\setlength{\parindent}{0pt}

\title{Driven Square Spin Ice as an Effective Nonequilibrium Vertex Model:\\
Assumptions, Coarse-Graining, and Steady-State Interpretation}
\author{}
\date{\today}

\begin{document}
\maketitle

\begin{abstract}
We present a detailed account of the effective model implemented in
\texttt{simulations/spinice/spinice.py}, with emphasis on the modeling
assumptions relevant to experts in spin ice, nonequilibrium statistical
mechanics, and periodically driven stochastic systems. The simulation should be
read as a coarse-grained, edge-based, finite-temperature vertex model for a
driven square artificial spin ice inspired by colloidal realizations. Its
degrees of freedom are binary occupancy variables on the lattice edges, its
static energetics are given by a local quadratic penalty on ice-rule
violations, and its dynamics are generated by a time-dependent Metropolis
kernel subject to a rotating external drive. The model therefore does not
represent microscopic colloidal trajectories, hydrodynamics, or continuous
barrier crossing. Rather, it provides a reduced description in which the
competition between local ice-rule energetics, thermal activation, and periodic
forcing produces a nonequilibrium steady state (NESS). We formulate the model
precisely, identify the implicit and explicit assumptions of the
coarse-graining, and clarify which conclusions are robust at the level of an
effective vertex theory and which would require a more microscopic treatment.
\end{abstract}

\section{Introduction}

Artificial spin-ice systems are routinely analyzed at several distinct levels
of description. At the most microscopic end, one may attempt to resolve the
elementary objects themselves, including their geometry, interactions, barrier
crossing events, and environment-induced dissipation. At the most coarse-grained
end, one may instead represent the system as a constrained statistical model
whose relevant variables are edge orientations and whose low-energy structure is
organized by vertex rules. The code studied here belongs firmly to the latter
category.

The central object of interest is a square-lattice ice manifold driven by a
time-dependent protocol that periodically biases horizontal and vertical edges
out of phase. The intended phenomenology is nonequilibrium: the drive competes
with the local energetic preference for 2-in-2-out vertices and with
thermally activated rearrangements, thereby generating a driven steady state
rather than an equilibrium Gibbs ensemble.

For an expert audience, the key question is not simply ``what does the code
do?'' but rather ``what effective stochastic process is actually being defined,
and under what assumptions can its output be interpreted as physically
meaningful?'' This note answers that question in detail.

\section{Degrees of Freedom and Geometric Structure}

The model is defined on an $L \times L$ square lattice with periodic boundary
conditions in both directions. The periodicity eliminates boundary effects and
places the system on a torus, so that all vertices are topologically
equivalent.

The elementary dynamical variables live on the \emph{edges} rather than on the
vertices. Each edge $e$ carries a binary variable
\[
  s_e \in \{-1,+1\}.
\]
The physical interpretation adopted by the code is that of a bistable trap
located on each edge: the sign of $s_e$ determines toward which endpoint vertex
the trapped particle is displaced. For horizontal edges, $s_e=+1$ means
displacement toward the right endpoint and $s_e=-1$ toward the left endpoint.
For vertical edges, $s_e=+1$ means displacement upward and $s_e=-1$ downward.

This is already the first nontrivial modeling decision. The code does not
resolve the continuous coordinate of a colloid inside a double well, nor does it
represent the full trap profile. It assumes from the outset that all
intra-trap dynamics may be compressed into a two-state variable. The resulting
model is therefore not a first-principles colloidal dynamics, but an effective
edge-Ising representation motivated by the colloidal picture.

There are $L^2$ vertices and $2L^2$ edge variables. The simulation precomputes
incidence tables that identify, for every vertex, its four incident edges and
the sign convention under which an edge counts as pointing into that vertex.

\section{Vertex Rule, Charge Variable, and Local Ice Energetics}

At each vertex $v$, the code counts the number of incident edges pointing
inward. Denoting this number by $n_v$, one introduces the local charge-like
deviation
\[
  q_v = n_v - 2.
\]
The 2-in-2-out ice manifold is therefore characterized by $q_v=0$ at every
vertex. Single defects correspond to $q_v=\pm1$ and double defects to
$q_v=\pm2$.

The static energy functional is taken to be
\[
  H_{\mathrm{vertex}}[s] = J \sum_v q_v^2
  = J\sum_v (n_v-2)^2.
\]
This choice encodes a standard hierarchy of defect costs:
\[
  E(q_v=0)=0,\qquad E(q_v=\pm1)=J,\qquad E(q_v=\pm2)=4J.
\]

Experts will immediately recognize the significance of this assumption. The
simulation is not employing a microscopic interaction model from which the
vertex spectrum is derived. Instead, it posits the vertex spectrum directly and
treats it as the relevant coarse-grained energetics. All 2-in-2-out states are
therefore exactly degenerate, and any energetic distinctions internal to the ice
manifold are discarded.

This approximation has several consequences.

First, the model is effectively local at the level of the Hamiltonian. Beyond
the fact that a single edge contributes to two neighboring vertices, there is no
explicit long-range coupling between defects. No dipolar tails, no Coulomb-like
entropic interactions, and no geometry-dependent subleading splittings are
included.

Second, the parameter $J$ should not be understood as a microscopic pair
interaction. It is an effective scale controlling the energetic suppression of
local ice-rule violations in the reduced model. Interpreting $J$ requires care:
it may phenomenologically encode several microscopic effects at once.

\section{Time-Dependent Drive and the Effective Single-Edge Tilt}

The nonequilibrium ingredient is a spatially uniform rotating drive
\[
  \mathbf{F}(t) = F_0 \bigl(\cos \omega t, \sin \omega t\bigr).
\]
For each edge $e$, the code defines a unit vector $\hat{\mathbf e}_e$ oriented
from the ``minus'' endpoint toward the ``plus'' endpoint. The instantaneous edge
bias is then
\[
  h_e(t) = \mathbf{F}(t)\cdot \hat{\mathbf e}_e.
\]
The drive contribution to the energy is
\[
  H_{\mathrm{drive}}[s,t] = - \sum_e h_e(t)s_e.
\]

For a horizontal edge,
\[
  h_e(t)=F_0\cos(\omega t),
\]
whereas for a vertical edge,
\[
  h_e(t)=F_0\sin(\omega t).
\]
Thus the horizontal and vertical sectors are driven with a quarter-period phase
shift. This is the sense in which the protocol is non-reciprocal in the present
effective description: the bias pattern felt by one orientational sector is not
synchronous with that felt by the orthogonal sector.

At the level of a single isolated edge, one may regard the drive term as the
coarse-grained remnant of a time-dependent tilt of an underlying bistable
potential. If one imagines a symmetric double well reduced to two minima, then
$h_e(t)$ provides the instantaneous asymmetry between the two wells. However,
the barrier itself is not explicitly modeled. The code retains only the
two-state energy splitting, not the full continuous landscape.

\section{Total Effective Hamiltonian}

The instantaneous energy function used by the dynamics is
\[
  H[s,t] = J\sum_v q_v^2 - \sum_e h_e(t)s_e.
\]
This Hamiltonian is explicitly time dependent. As such, there is no stationary
Boltzmann distribution of the form $\propto e^{-\beta H[s]}$ associated with
the full long-time process. The relevant object is instead the stochastic
dynamics induced by this time-dependent energy landscape.

It is useful to emphasize that the Hamiltonian has two conceptually distinct
roles:

\begin{enumerate}
\item the vertex term defines the static interaction structure and the local
energetic hierarchy of ice-rule compliance versus defect creation;
\item the drive term provides a time-periodic conjugate field acting directly on
the edge variables.
\end{enumerate}

The nonequilibrium response is generated by the competition of these two
contributions under stochastic updating.

\section{Stochastic Dynamics}

\subsection{Metropolis dynamics as an effective kinetic rule}

The code evolves the edge variables by single-edge flip attempts
$s_k \to -s_k$. A proposed flip changes the energy by
\[
  \Delta E = \Delta E_{\mathrm{vertex}} + \Delta E_{\mathrm{drive}}.
\]
Because the vertex term is local, only the two endpoint vertices of edge $k$
need to be updated. The drive contribution is simply the change in the
time-dependent single-edge bias.

The acceptance probability is the standard Metropolis rule
\[
  P_{\mathrm{acc}} = \min\{1,e^{-\beta \Delta E}\},
  \qquad \beta = 1/T.
\]
One Monte Carlo sweep consists of $N_e=2L^2$ randomly selected flip attempts.

This is the second major modeling assumption. The code does not derive
transition rates from barrier heights or from an overdamped Langevin dynamics in
continuous space. Instead, it adopts Metropolis updating as an effective kinetic
prescription. Temperature therefore enters only through the stochastic
acceptance of energetically unfavorable moves.

For equilibrium spin systems this rule is standard and unambiguous. In the
present driven setting, it should be interpreted more cautiously: it defines a
time-dependent Markov chain whose instantaneous local bias is Boltzmann-like, but
whose global long-time statistics are not equilibrium statistics.

\subsection{Monte Carlo time and discretization of the drive}

The simulation introduces a Monte Carlo time variable $t_{\mathrm{MC}}$. After
each sweep, time advances by one unit,
\[
  t_{\mathrm{MC}} \mapsto t_{\mathrm{MC}}+1.
\]
The drive $\mathbf{F}(t)$ is evaluated once per sweep and held fixed throughout
the $N_e$ flip attempts of that sweep.

This is a nontrivial time discretization. It implies that the field varies on
the scale of sweeps, not on the scale of individual elementary updates. In other
words, the piecewise-constant approximation to the drive is built directly into
the dynamics. If $\omega$ becomes too large relative to the sweep time, the
distinction between an ideal continuously rotating drive and the implemented
piecewise-constant protocol becomes quantitatively important.

\subsection{Temperature as effective thermal activation}

The parameter $T$ in the code is an effective Monte Carlo temperature. It should
not automatically be identified with a thermodynamic bath temperature in a
microscopic colloidal fluid model. What the code assumes is simply that the
ratio of upward to downward local moves is controlled by the Boltzmann factor
$e^{-\beta \Delta E}$ computed from the coarse-grained energy.

This is entirely reasonable as an effective statistical mechanics ansatz, but it
places the simulation in the class of phenomenological driven stochastic lattice
models rather than in the class of explicit solvent-coupled particle dynamics.

\section{Nonequilibrium Steady State and Detailed Balance}

Without driving, the Metropolis kernel would sample the equilibrium ensemble
associated with $H_{\mathrm{vertex}}$, modulo the usual assumptions of
ergodicity and equilibration. With the rotating drive present, the situation is
qualitatively different. The energy function changes in time, and therefore the
stochastic process no longer obeys detailed balance with respect to any static
measure.

The appropriate asymptotic object is a periodically driven nonequilibrium steady
state. More precisely, after transients decay, the probability distribution over
configurations is expected to become time-periodic with the same period as the
drive, or at least to exhibit stationary cycle-averaged statistics. The code
then samples observables over a long measurement window and reports time
averages such as the mean defect density.

From the standpoint of nonequilibrium statistical mechanics, the simulation is
best interpreted as defining a periodically forced Markov chain on the
configuration space of square-ice edge variables. The observables extracted from
the run are empirical estimates of the cycle-averaged response of this chain.

\section{Warmup, Measurement, and the Operational Definition of NESS}

The script divides the run into two stages. During a warmup interval of
\texttt{n\_warmup} sweeps, the drive is already active and no data are retained.
During a measurement interval of \texttt{n\_measure} sweeps, observables are
sampled after each sweep and averaged.

Operationally, the NESS is therefore defined numerically rather than
theoretically: one assumes that the warmup interval is sufficient to remove the
dependence on the initial condition and that the subsequent sampling window is
long enough to characterize the long-time driven regime.

This procedure is standard, but experts will appreciate the associated caveats.
If relaxation is slow, if there are long-lived metastable sectors, or if the
drive induces broad temporal correlations, then finite warmup and finite
measurement windows may bias the results. In particular, quoted error bars based
on the naive standard deviation over sweeps are not equivalent to a full
autocorrelation-aware uncertainty analysis.

\section{Observables}

The primary observable is the defect density
\[
  \rho_{\mathrm{def}} =
  \frac{1}{L^2}\sum_v \mathbf{1}_{q_v \neq 0},
\]
namely the fraction of vertices that violate the local ice rule. This is a
natural measure of how far the driven system is pushed away from the ice
manifold.

The code also computes the mean squared charge density
\[
  \langle q^2 \rangle =
  \frac{1}{L^2}\sum_v q_v^2,
\]
which resolves the distinction between singly and doubly charged defects more
sensitively than $\rho_{\mathrm{def}}$ alone.

These observables are then used to construct line cuts in frequency and
two-parameter response maps in the $(\omega,F_0)$ plane. In the language of the
code, these maps are called phase diagrams, though for precision one should
usually describe them as nonequilibrium response diagrams or defect landscapes.
Unless singular thermodynamic behavior is independently established, one should
not infer sharp phase boundaries merely from smooth crossovers in these plots.

\section{Modeling Assumptions: Explicit and Implicit}

For clarity, it is worth gathering the assumptions in one place.

\subsection*{1. Two-state reduction}
Each edge trap is replaced by a binary variable. Continuous motion within the
trap and explicit barrier geometry are discarded.

\subsection*{2. Effective local energetics}
The entire static interaction structure is compressed into the local quadratic
vertex penalty $J\sum_v q_v^2$.

\subsection*{3. Degenerate ice manifold}
All 2-in-2-out states are exactly degenerate. Any internal energetic structure
within the ice manifold is neglected.

\subsection*{4. No explicit long-range forces}
There are no dipolar, Coulombic, hydrodynamic, or elastic interactions beyond
those encoded phenomenologically by the local vertex term.

\subsection*{5. Uniform periodic drive}
The forcing is perfectly sinusoidal, spatially uniform, and noiseless.
Disorder in amplitude, phase, or trap response is absent.

\subsection*{6. Effective thermal bath}
Temperature enters only through Metropolis acceptance. The model assumes that
the coarse-grained transitions are well represented by this effective stochastic
rule.

\subsection*{7. Sweep-level time discretization}
The rotating field is updated once per sweep, not continuously during the sweep.

\subsection*{8. Periodic boundaries}
Boundary effects, edge nucleation, and corner-specific defect physics are
deliberately excluded.

\subsection*{9. Finite-size and finite-time estimation}
Reported observables are finite-size, finite-time empirical averages over a
numerically defined driven steady state.

\section{What the Model Captures Reliably}

Within its intended scope, the model is well adapted to several questions of
genuine interest:

\begin{itemize}
\item how a rotating external bias competes with the local ice-rule constraint,
\item how defect production depends on drive amplitude and frequency,
\item how thermal activation modifies the driven response,
\item how a periodically forced constrained lattice system settles into a
nonequilibrium steady regime,
\item how one may visualize the structure of the resulting defect landscape in
parameter space.
\end{itemize}

In particular, the model is quite suitable for asking whether there exists a
frequency window in which the driven production of defects is enhanced or
suppressed, and how that window shifts with temperature and forcing amplitude.

\section{What the Model Does Not Capture}

The limitations are equally important.

The model does not describe explicit colloidal mechanics. There are no real
particle trajectories, no mobility matrix, no viscous drag law, no hydrodynamic
coupling, and no continuous barrier crossing. The word ``colloidal'' therefore
enters primarily at the level of physical interpretation of the two-state edge
variable, not at the level of the simulated equations of motion.

Nor does the model include long-range defect interactions or entropic Coulomb
physics known from more refined spin-ice descriptions. Defects are local charges
only in the sense of their classification and local energetic weight, not in the
sense of an emergent gauge-mediated interaction extending across the sample.

Finally, Monte Carlo time is not automatically physical time. Although the code
interprets $\omega$ as a frequency measured in inverse sweep units, any mapping
from sweeps to laboratory time would require additional kinetic assumptions not
present in the implementation.

\section{Interpretive Status of the Simulation}

The correct scientific status of the code is, in our view, the following. It is
an effective driven stochastic vertex model inspired by colloidal square spin
ice. Its degrees of freedom are edge occupancies, its static energy penalizes
ice-rule violations, and its dynamics are governed by a time-dependent
Metropolis rule in a rotating field. It should therefore be understood as a
minimal model of NESS formation in a constrained lattice system, rather than as
a microscopic simulation of a particular experimental platform.

This distinction is not a weakness; on the contrary, it is precisely what makes
the model analytically legible. Because the coarse-graining is strong and the
assumptions are transparent, one can attribute observed behavior directly to the
competition among a small number of well-controlled ingredients: local defect
cost, thermal stochasticity, and periodic bias.

\section{Possible Extensions}

For future work, several natural extensions suggest themselves.

One may introduce explicit barrier heights and transition-state rates rather than
a purely Metropolis rule. One may resolve the continuous coordinate of each
colloid and replace the two-state reduction with an overdamped Langevin
dynamics. One may include quenched disorder in trap depths or local thresholds,
which is often essential in artificial-ice experiments. One may also introduce
longer-range defect interactions or subleading energetic splittings among
ice-rule states.

From the NESS perspective, further observables would also be valuable: cycle
resolved hysteresis, work and entropy production, frequency-resolved response
functions, and two-time correlation or persistence diagnostics. These would
bring the model into closer contact with the broader language of driven
stochastic thermodynamics.

\section{Conclusion}

The model implemented in \texttt{spinice.py} is best read as a periodically
driven, finite-temperature, effective square-ice vertex model. It is built from
binary edge variables, a local quadratic penalty on ice-rule defects, and a
rotating external field that biases horizontal and vertical edges out of phase.
Its stochastic dynamics are defined by a time-dependent Metropolis kernel, so
the long-time regime is a nonequilibrium steady state rather than an equilibrium
ensemble.

For experts, the crucial point is that the simulation is simultaneously modest
and useful. It is modest because it omits microscopic colloidal dynamics,
hydrodynamics, long-range interactions, and continuous-time kinetics. It is
useful because it isolates, in the clearest possible form, how a constrained
spin-ice manifold responds to thermal noise and periodic forcing. As a result,
the code provides a sensible starting point for qualitative and semi-quantitative
investigations of defect production and NESS structure in driven square spin
ice, provided its coarse-grained status is kept firmly in view.

\end{document}
